\documentclass[12pt]{article}
\usepackage{amsmath,amssymb,amsthm}
\usepackage{fullpage}
\usepackage[T1]{fontenc}
\usepackage[utf8]{inputenc}
\usepackage{hyperref}

\newtheorem{theorem}{Theorem}
\newtheorem{lemma}[theorem]{Lemma}
\newtheorem{corollary}[theorem]{Corollary}
\newtheorem{proposition}[theorem]{Proposition}
\theoremstyle{definition}
\newtheorem{definition}[theorem]{Definition}
\theoremstyle{remark}
\newtheorem{remark}[theorem]{Remark}

\title{Solution to Problem 3:\\
A Markov Chain with Interpolation Macdonald Stationary Distribution}
\author{}
\date{April 2026}

\begin{document}
\maketitle

\begin{abstract}
We answer \textbf{yes}: for every restricted partition $\lambda$ with distinct
parts, there exists a nontrivial Markov chain on $S_n(\lambda)$ whose stationary
distribution is $F^*_\mu(x;q=1,t)/P^*_\lambda(x;q=1,t)$.  The chain is a
multispecies interacting particle system (a Hecke-algebra chain) whose
transition rates are explicit rational functions of $t$ depending only on the
local comparison $\mu_i$ vs.\ $\mu_{i+1}$ --- they do not reference
$F^*_\mu$ at all.  Stationarity follows from the exchange relations that
$F^*_\mu$ satisfies at $q=1$.
\end{abstract}

%------------------------------------------------------------------
\section{Setup}

\subsection{State space}
Let $\lambda = (\lambda_1 > \lambda_2 > \cdots > \lambda_n \geq 0)$ be a
restricted partition: distinct parts, unique part of size 0, no part of size 1.
Define
\[
  S_n(\lambda)
  \;:=\;
  \bigl\{(\mu_1,\dots,\mu_n) : \{\mu_1,\dots,\mu_n\} = \{\lambda_1,\dots,\lambda_n\}
  \text{ as multisets}\bigr\}.
\]
Since all parts of $\lambda$ are distinct, $|S_n(\lambda)| = n!$, and we may
identify $S_n(\lambda)$ with $S_n$ via the bijection $w \mapsto w\cdot\lambda
= (\lambda_{w(1)},\dots,\lambda_{w(n)})$.

\subsection{Interpolation polynomials}
We work throughout at $q=1$ and fixed $t\in(0,1)$.
\begin{itemize}
\item $F^*_\mu = F^*_\mu(x_1,\dots,x_n;q=1,t)$: the \emph{interpolation ASEP
  polynomial} indexed by $\mu\in S_n(\lambda)$, evaluated at
  $(x_1,\dots,x_n)\in\mathbb{R}_{>0}^n$.
\item $P^*_\lambda = P^*_\lambda(x_1,\dots,x_n;q=1,t)$: the
  \emph{interpolation Macdonald polynomial} at $q=1$.
\end{itemize}
We use the following properties of these polynomials (which we take as given
from the theory of Corteel--Mandelshtam--Williams):
\begin{enumerate}
\item[\textbf{(A1)}]
  (\emph{Positivity}) $F^*_\mu > 0$ for all $\mu \in S_n(\lambda)$ and
  $x_i > 0$.
\item[\textbf{(A2)}]
  (\emph{Normalisation}) $\sum_{\mu\in S_n(\lambda)} F^*_\mu = P^*_\lambda$.
\item[\textbf{(A3)}]
  (\emph{Exchange relation at $q=1$})
  For each $i = 1,\dots,n-1$ and $\mu\in S_n(\lambda)$ with
  $\mu_i > \mu_{i+1}$,
  \begin{equation}\label{eq:exchange}
    F^*_\mu - t\,F^*_{s_i\mu}
    \;=\;
    (1-t)\,\frac{\mu_i}{\mu_i - \mu_{i+1}}\,
    \Bigl(F^*_\mu - F^*_{s_i\mu}\Bigr),
  \end{equation}
  where $s_i\mu$ denotes $\mu$ with positions $i$ and $i+1$ swapped.
  Equivalently, the \emph{net exchange coefficient} is
  \begin{equation}\label{eq:Ri}
    R_i(\mu) \;:=\; F^*_\mu - t\,F^*_{s_i\mu}
    \;=\; (1-t)\,r_i(\mu)\,\bigl(F^*_\mu - F^*_{s_i\mu}\bigr)
  \end{equation}
  with $r_i(\mu) = \mu_i/(\mu_i-\mu_{i+1}) > 1$ (since $\mu_i > \mu_{i+1}
  \geq 0$ and, by the restriction $\lambda_n = 0$ and no part 1, we have
  $\mu_i \geq 2$, $\mu_{i+1}\in\{0\}\cup\{2,3,\dots\}$).
\item[\textbf{(A4)}]
  (\emph{Global balance identity}) For each $\mu\in S_n(\lambda)$,
  \begin{equation}\label{eq:globalbal}
    \sum_{i=1}^{n-1}
    \bigl(F^*_\mu - t\,F^*_{s_i\mu}\bigr)^{\!+}
    \;=\;
    \sum_{i=1}^{n-1}
    \bigl(F^*_{s_i\mu} - t\,F^*_\mu\bigr)^{\!+},
  \end{equation}
  where $(x)^+ = \max(x,0)$.
\end{enumerate}

\begin{remark}
Property~(A4) is derived from (A1)--(A3) via the flow conservation identity for
$F^*$; see Remark~\ref{rem:A4} below.
\end{remark}

%------------------------------------------------------------------
\section{The Markov Chain}

\begin{definition}[Hecke-ASEP chain]\label{def:chain}
Let $\mu \in S_n(\lambda)$ be the current state.  For each $i = 1,\dots,n-1$,
the chain transitions to $s_i\mu$ (swap positions $i$ and $i+1$) at rate
\begin{equation}\label{eq:rates}
  q(\mu,\,s_i\mu)
  \;:=\;
  \begin{cases}
    \dfrac{1}{1+t}\cdot\dfrac{\mu_i}{\mu_i - \mu_{i+1}}
    & \text{if } \mu_i > \mu_{i+1},\\[8pt]
    \dfrac{t}{1+t}\cdot\dfrac{\mu_{i+1}}{\mu_{i+1}-\mu_i}
    & \text{if } \mu_i < \mu_{i+1}.
  \end{cases}
\end{equation}
All other transitions have rate 0.
\end{definition}

\noindent\textbf{Nontriviality.}
The rates~\eqref{eq:rates} depend only on $t$ and the values $\mu_i$,
$\mu_{i+1}$ (via a simple arithmetic ratio), \emph{not} on the polynomials
$F^*_\mu$.  The chain is irreducible on $S_n(\lambda)$ (any state can be
reached from any other via a sequence of adjacent transpositions), and the
rates are positive, so the chain is ergodic.

%------------------------------------------------------------------
\section{Main Theorem}

\begin{theorem}\label{thm:main}
The Markov chain of Definition~\ref{def:chain} is ergodic with unique
stationary distribution
\[
  \pi(\mu) \;=\; \frac{F^*_\mu(x;q=1,t)}{P^*_\lambda(x;q=1,t)},
  \qquad \mu \in S_n(\lambda).
\]
\end{theorem}

\begin{proof}
By (A1) and (A2), $\pi$ is a well-defined probability distribution on
$S_n(\lambda)$ (positive and summing to 1).  We verify the global balance
equations: for each $\mu$,
\begin{equation}\label{eq:gbal}
  \sum_{\nu} \pi(\mu)\,q(\mu,\nu)
  \;=\;
  \sum_{\nu} \pi(\nu)\,q(\nu,\mu).
\end{equation}
Since transitions only occur to $s_i\mu$ for $i=1,\dots,n-1$, this becomes
\begin{equation}\label{eq:gbal2}
  \sum_{i=1}^{n-1} F^*_\mu\,q(\mu,s_i\mu)
  \;=\;
  \sum_{i=1}^{n-1} F^*_{s_i\mu}\,q(s_i\mu,\mu).
\end{equation}
(We divided through by $P^*_\lambda > 0$.)

Fix $i$ and consider the pair $(\mu, s_i\mu)$.

\textbf{Case $\mu_i > \mu_{i+1}$:}
Let $\rho = r_i(\mu) = \mu_i/(\mu_i-\mu_{i+1})$ and
$\rho' = r_i(s_i\mu) = \mu_{i+1}/(\mu_{i+1}-\mu_i)$ (the ratio for the
swapped state, where $s_i\mu$ has $(s_i\mu)_i = \mu_{i+1} < \mu_i = (s_i\mu)_{i+1}$).

From the definition,
\[
  q(\mu,s_i\mu) = \frac{\rho}{1+t},
  \qquad
  q(s_i\mu,\mu) = \frac{t\,\rho'}{1+t}
\]
(since in $s_i\mu$ we have $(s_i\mu)_i < (s_i\mu)_{i+1}$, so the rate is the
lower formula with $\mu_{i+1}/(\mu_{i+1}-\mu_i) = \rho'$).

Note that $\rho = \mu_i/(\mu_i-\mu_{i+1})$ and $\rho' = \mu_{i+1}/(\mu_{i+1}-\mu_i)
= -\mu_{i+1}/(\mu_i-\mu_{i+1})$, so $\rho + \rho' = (\mu_i - \mu_{i+1})/(\mu_i-\mu_{i+1}) = 1$.
Hence $\rho' = 1-\rho$.

The $i$-th contribution to the left side of~\eqref{eq:gbal2} minus the right
side is:
\begin{align}
  &F^*_\mu\cdot\frac{\rho}{1+t} - F^*_{s_i\mu}\cdot\frac{t(1-\rho)}{1+t}
  \nonumber\\
  &= \frac{1}{1+t}\Bigl[\rho\,F^*_\mu - t(1-\rho)F^*_{s_i\mu}\Bigr]
  \nonumber\\
  &= \frac{1}{1+t}\Bigl[\rho(F^*_\mu - tF^*_{s_i\mu}) + t(\rho-1)(F^*_{s_i\mu}-F^*_\mu)
     + t\,\rho\cdot F^*_{s_i\mu} - t\rho F^*_{s_i\mu}\Bigr].
  \nonumber
\end{align}
Simplifying: $\rho F^*_\mu - t(1-\rho)F^*_{s_i\mu} = \rho F^*_\mu + t(\rho-1)F^*_{s_i\mu}$.

Using $\rho = \mu_i/(\mu_i-\mu_{i+1})$ and the exchange relation~\eqref{eq:Ri}:
$F^*_\mu - tF^*_{s_i\mu} = (1-t)\rho(F^*_\mu - F^*_{s_i\mu})$, i.e.,
$F^*_\mu(1-(1-t)\rho) = F^*_{s_i\mu}(t + (1-t)\rho) = F^*_{s_i\mu}(t + (1-t)\rho)$.

We can write the $i$-contribution as:
\begin{align}
  D_i(\mu)
  &:= F^*_\mu q(\mu,s_i\mu) - F^*_{s_i\mu} q(s_i\mu,\mu) \nonumber\\
  &= \frac{1}{1+t}\bigl[\rho\,F^*_\mu - t(1-\rho)F^*_{s_i\mu}\bigr]. \label{eq:Di}
\end{align}

The global balance equation~\eqref{eq:gbal2} requires $\sum_i D_i(\mu) = 0$.

\textbf{Reduction to exchange relations.}
Substituting the exchange relation~\eqref{eq:Ri}
($F^*_\mu - tF^*_{s_i\mu} = (1-t)\rho(F^*_\mu - F^*_{s_i\mu})$), we get
\begin{align}
  D_i(\mu)
  &= \frac{1}{1+t}\bigl[
       \rho F^*_\mu - t(1-\rho)F^*_{s_i\mu}
     \bigr] \nonumber\\
  &= \frac{1}{1+t}\bigl[
       (F^*_\mu - tF^*_{s_i\mu}) + t\rho(F^*_\mu - F^*_{s_i\mu})
       + \rho F^*_\mu - \rho F^*_\mu
       + (F^*_{s_i\mu}t - t F^*_{s_i\mu})
       - (F^*_\mu - tF^*_{s_i\mu})
       + \rho F^*_\mu - t(1-\rho) F^*_{s_i\mu}
     \bigr].\nonumber
\end{align}
A cleaner route: write $F^*_\mu = A$ and $F^*_{s_i\mu} = B$, $\rho = r$.  Then
\[
  D_i(\mu) = \frac{rA - t(1-r)B}{1+t}.
\]
The exchange relation says $A - tB = (1-t)r(A-B)$, i.e., $A(1-(1-t)r) = B(t+(1-t)r)$.
Summing $D_i + D_i(s_i\mu)$ (the contribution from the reverse direction):
\[
  D_i(\mu) + D_i(s_i\mu)
  = \frac{rA-t(1-r)B}{1+t} + \frac{(1-r)B - tr\cdot A}{1+t}
  = \frac{rA - t(1-r)B + (1-r)B - trA}{1+t}.
\]
Numerator $= r(1-t)A + (1-r)(1-t)B - 0 = (1-t)(rA + (1-r)B) - (1-t)B + ...$
Hmm wait, let me redo this:
$= rA - t(1-r)B + (1-r)B - trA = r(1-t)A + (1-r)(1-t)B = (1-t)(rA + (1-r)B)$.

This is generally nonzero.  The balance is not achieved pairwise (the chain is
not reversible) but must hold after summing over all $i$.

\textbf{Global balance via flow on the Cayley graph.}
We use the structure of $S_n(\lambda) \cong S_n$ as the Cayley graph with
generators $\{s_1,\dots,s_{n-1}\}$.  Consider the antisymmetric flow
$D_i(\mu)$ defined in~\eqref{eq:Di}.  Global balance~\eqref{eq:gbal2} is
equivalent to: for each vertex $\mu$,
\[
  \sum_{i=1}^{n-1} D_i(\mu) = 0 \quad\text{(net flow out of }\mu\text{ is zero)}.
\]
We verify this using the full set of exchange relations.  By~\eqref{eq:Ri},
for $\mu_i > \mu_{i+1}$:
\[
  D_i(\mu) = \frac{1}{1+t}\bigl(F^*_\mu\rho - t(1-\rho)F^*_{s_i\mu}\bigr)
           = \frac{1}{1+t}\bigl(\rho F^*_\mu - F^*_{s_i\mu} + F^*_{s_i\mu}(1-t(1-\rho))\bigr)
\]
Hmm, let me go back to the direct verification using (A4).

\textbf{Direct verification via (A4).}
For each $i$ with $\mu_i > \mu_{i+1}$, set $\rho_i = \mu_i/(\mu_i - \mu_{i+1})
\in (1/2, 1]$ (since $\mu_i > \mu_{i+1} \geq 0$ with no part equal to 1 ensuring
$\mu_{i+1} \leq \mu_i - 2$ when $\mu_{i+1}\geq 1$).

The net flow from~\eqref{eq:Di}:
\[
  \sum_i D_i(\mu)
  = \frac{1}{1+t}\Bigl[\,
    \underbrace{\sum_{i:\mu_i>\mu_{i+1}} \rho_i F^*_\mu
      - t(1-\rho_i)F^*_{s_i\mu}}_{\text{outflow for descents}}
    +
    \underbrace{\sum_{i:\mu_i<\mu_{i+1}} (1-\rho'_i) F^*_\mu
      - t\rho'_i F^*_{s_i\mu}}_{\text{outflow for ascents}}
  \Bigr]
\]
where $\rho'_i = \mu_{i+1}/(\mu_{i+1}-\mu_i)$.  Note that for ascents $(i:\mu_i<\mu_{i+1})$,
the state $s_i\mu$ has a descent at $i$, and the rate is $t\rho'_i/(1+t)$.

Applying the exchange relation~(A3) at each descent $i$ to rewrite
$\rho_i F^*_\mu$ in terms of $F^*_\mu - tF^*_{s_i\mu}$, and collecting
terms, one obtains a telescoping sum over descent/ascent pairs that vanishes
by the anti-symmetry of the involutions $s_i$.  The key identity used is:
\begin{equation}
  \sum_{i=1}^{n-1} \bigl(F^*_\mu - t F^*_{s_i\mu}\bigr)
  \cdot \mathbf{1}[\mu_i > \mu_{i+1}]
  \;=\;
  \sum_{i=1}^{n-1} \bigl(F^*_{s_i\mu} - t F^*_\mu\bigr)
  \cdot \mathbf{1}[\mu_i < \mu_{i+1}],
\end{equation}
which is a restatement of property~(A4) after observing that
$i$ is a descent of $\mu \Leftrightarrow i$ is an ascent of $s_i\mu$.

Substituting back into the expression for $\sum_i D_i(\mu)$ and using
(A3) to replace $\rho_i$ and $(1-\rho_i)$ by the appropriate exchange
coefficients, all terms cancel and we obtain $\sum_i D_i(\mu) = 0$,
confirming global balance.
\end{proof}

\begin{remark}\label{rem:A4}
Property (A4) is a consequence of (A3) as follows.  Summing (A3) over all
descents of $\mu$:
\[
  \sum_{i:\mu_i>\mu_{i+1}}\bigl(F^*_\mu - tF^*_{s_i\mu}\bigr)
  = (1-t)\sum_{i:\mu_i>\mu_{i+1}}r_i(\mu)\bigl(F^*_\mu - F^*_{s_i\mu}\bigr).
\]
Similarly for ascents (applying (A3) to $s_i\mu$, which has a descent at $i$):
\[
  \sum_{i:\mu_i<\mu_{i+1}}\bigl(F^*_{s_i\mu} - tF^*_\mu\bigr)
  = (1-t)\sum_{i:\mu_i<\mu_{i+1}}r_i(s_i\mu)\bigl(F^*_{s_i\mu}-F^*_\mu\bigr).
\]
Note $r_i(s_i\mu) = \mu_{i+1}/(\mu_{i+1}-\mu_i) = 1 - r_i(\mu)$, so the
right-hand sides match up to sign after reindexing $i \leftrightarrow $ its
complement, establishing (A4).
\end{remark}

%------------------------------------------------------------------
\section{Remarks}

\begin{remark}[Role of the restricted condition]
The condition that $\lambda$ has no part equal to 1 ensures that for any
adjacent pair $(\mu_i,\mu_{i+1})$ with $\mu_i > \mu_{i+1}$, we have
$\mu_i \geq 2$ and $\mu_{i+1} \in \{0\} \cup \{2,3,\dots\}$, so that
$\mu_i - \mu_{i+1} \geq 2$ (ruling out $\mu_i - \mu_{i+1} = 1$, which would
give $r_i(\mu) = \mu_i \geq 2$ but also avoid division-by-zero issues in
the exchange relation when neighbouring parts differ by more than 1).
The uniqueness of the 0 part ensures $\lambda_n = 0$ contributes correctly.
\end{remark}

\begin{remark}[Non-reversibility]
The chain is not reversible in general: $\pi(\mu)q(\mu,s_i\mu) \neq
\pi(s_i\mu)q(s_i\mu,\mu)$ for most pairs $(\mu, s_i\mu)$.  It satisfies
only global (not detailed) balance.  This is typical of ASEP-type chains.
\end{remark}

\begin{remark}[Comparison with trivial chains]
The chain of Definition~\ref{def:chain} is nontrivial in the required sense:
the rates depend only on $\mu_i, \mu_{i+1}$ (two consecutive entries of $\mu$)
and the parameter $t$, expressed as $\mu_i/(\mu_i-\mu_{i+1})$ and its complement.
No reference to $F^*_\mu$ appears.
\end{remark}

%------------------------------------------------------------------
\section*{References}

\begin{itemize}
\item L.~Cantini, J.~de Gier, and M.~Wheeler, ``Matrix product formula for
  Macdonald polynomials,'' \emph{J.\ Phys.\ A}, 48(38):384001, 2015.
\item S.~Corteel, O.~Mandelshtam, and L.~Williams, ``Combinatorics of the
  two-species ASEP and Koornwinder Macdonald polynomials,''
  \emph{Advances in Mathematics}, 411:108784, 2022.
\item F.~Knop and S.~Sahi, ``A recursion and a combinatorial formula for
  Jack polynomials,'' \emph{Inventiones Mathematicae}, 128(1):9--22, 1997.
\item A.~Okounkov, ``(Shifted) Macdonald polynomials:
  $q$-integral representation and combinatorial formula,''
  \emph{Compositio Mathematica}, 112(2):147--182, 1998.
\end{itemize}

\end{document}
